\epigraph{The paper analyzes investment policy for a firm which cannot disinvest faster than its capital depreciates.}{\textit{Kenneth Arrow}, ``Optimal Capital Policy with Irreversible Investment'' (1968)}

\epigraph{This paper determines the time series behavior of investment, output, and prices in a competitive industry with a stochastic demand.}{\textit{Robert Lucas and Edward Prescott}, ``Investment Under Uncertainty'' (1971)}



% ==============================================================
\section{Risk and Uncertainty}
\label{sec:ac-risk}
% ==============================================================

We now introduce uncertainty into the adjustment cost framework. This section connects three strands: adjustment costs (the CES transformation), portfolio choice (allocation across risky capital types), and consumption-savings under uncertainty. The key insight is that the timing of information---whether shocks are realized before or after investment decisions---fundamentally changes the problem's structure and maps to different CES aggregations.

% --------------------------------------------------------------
\subsection{Sources of Uncertainty}
\label{subsec:ac-risk-sources}
% --------------------------------------------------------------

Consider an economy with $N$ capital types and $S$ states of the world. State $s$ occurs with probability $\pi_s > 0$, where $\sum_{s=1}^{S} \pi_s = 1$.

\paragraph{Capital-Specific Productivity Shocks.}
The productivity of capital type $i$ in state $s$ is $a_{is}$. Output in state $s$ is:
\begin{equation}
Y_s = \F(\bk, \ba_s) = \left( \sum_{i=1}^{N} (\alpha_i a_{is})^{1-\sigma} k_i^{\sigma} \right)^{1/\sigma},
\label{eq:ac-risk-production}
\end{equation}
where $\ba_s = (a_{1s}, \ldots, a_{Ns})$ is the vector of productivity realizations in state $s$.

\paragraph{Two Timing Conventions.}
The crucial modeling choice is \emph{when} uncertainty resolves relative to \emph{when} decisions are made:

\begin{enumerate}[label=(\roman*)]
\item \textbf{Ex-post uncertainty}: Investment $\bI$ is chosen \emph{before} the state $s$ is revealed. The household faces a portfolio problem: allocate investment across capital types not knowing which will be productive.

\item \textbf{Ex-ante uncertainty}: Investment $\bI_s$ is chosen \emph{after} the state $s$ is revealed (or equivalently, state-contingent investment is possible). The household can tailor investment to the realized state.
\end{enumerate}

These two cases lead to different aggregation structures and different roles for the CES machinery.

% --------------------------------------------------------------
\subsection{Ex-Post Uncertainty: The Portfolio Problem}
\label{subsec:ac-risk-expost}
% --------------------------------------------------------------

Suppose investment is chosen before uncertainty resolves. This is the natural timing for many macroeconomic applications: firms commit to investment plans, and productivity is realized later.

\paragraph{Setup.}
The household chooses $(C_0, \bI)$ in period 0, then the state $s$ is revealed, and period-1 consumption is:
\begin{equation}
C_{1s} = \F(\bI, \ba_s) = \left( \sum_{i=1}^{N} (\alpha_i a_{is})^{1-\sigma} I_i^{\sigma} \right)^{1/\sigma}.
\label{eq:ac-risk-C1s}
\end{equation}

The transformation constraint applies in period 0 (before uncertainty):
\begin{equation}
\Hfun(C_0, \bI) = Y_0.
\label{eq:ac-risk-H0}
\end{equation}

\paragraph{Preferences.}
With expected utility:
\begin{equation}
U = u(C_0) + \beta \sum_{s=1}^{S} \pi_s \, u(C_{1s}) = u(C_0) + \beta \, \mathbb{E}[u(C_1)].
\label{eq:ac-risk-EU}
\end{equation}

With Epstein-Zin preferences (separating risk aversion $\gamma$ from IES $\theta$):
\begin{equation}
U = u(C_0) + \beta \, u\left( \CE(C_1) \right),
\label{eq:ac-risk-EZ}
\end{equation}
where the certainty equivalent is:
\begin{equation}
\CE(C_1) = \left( \sum_{s=1}^{S} \pi_s \, C_{1s}^{1-\gamma} \right)^{\frac{1}{1-\gamma}}.
\label{eq:ac-risk-CE}
\end{equation}

\paragraph{The Nested Problem.}
The problem has a nested structure:
\begin{enumerate}[nosep]
\item \textbf{Inner problem (portfolio)}: Given total investment $\bar{I}$, choose portfolio shares $\boldsymbol{\phi}$ to maximize the certainty equivalent of future consumption.
\item \textbf{Outer problem (consumption-savings)}: Choose $C_0$ and $\bar{I}$ given the optimal certainty equivalent.
\end{enumerate}

\paragraph{The Portfolio Problem.}
Let $\phi_i = I_i / \bar{I}$ with $\sum_i \phi_i = 1$. Future consumption in state $s$ is:
\begin{equation}
C_{1s} = \bar{I} \cdot \F(\boldsymbol{\phi}, \ba_s).
\label{eq:ac-risk-C1s-phi}
\end{equation}

The certainty equivalent becomes:
\begin{equation}
\CE(C_1) = \bar{I} \cdot \left( \sum_{s=1}^{S} \pi_s \, \F(\boldsymbol{\phi}, \ba_s)^{1-\gamma} \right)^{\frac{1}{1-\gamma}} \equiv \bar{I} \cdot \mathcal{R}(\boldsymbol{\phi}),
\label{eq:ac-risk-CE-phi}
\end{equation}
where $\mathcal{R}(\boldsymbol{\phi})$ is the \emph{risk-adjusted return} per unit of investment:
\begin{equation}
\mathcal{R}(\boldsymbol{\phi}) = \CE_s\left[ \F(\boldsymbol{\phi}, \ba_s) \right] = \left( \sum_{s=1}^{S} \pi_s \, \F(\boldsymbol{\phi}, \ba_s)^{1-\gamma} \right)^{\frac{1}{1-\gamma}}.
\label{eq:ac-risk-R-phi}
\end{equation}

The portfolio problem is:
\begin{equation}
\max_{\boldsymbol{\phi}} \; \mathcal{R}(\boldsymbol{\phi}) \quad \text{subject to} \quad \sum_{i=1}^{N} \phi_i = 1.
\label{eq:ac-risk-portfolio}
\end{equation}

\begin{calloutimportant}[A CES of CES Problem]
The objective $\mathcal{R}(\boldsymbol{\phi})$ is a \emph{certainty equivalent} (CES with power $1-\gamma$ over states) of a \emph{production function} (CES with power $\sigma$ over capital types):
\[
\mathcal{R}(\boldsymbol{\phi}) = \underbrace{\left( \sum_{s} \pi_s \, \left[ \underbrace{\left( \sum_{i} (\alpha_i a_{is})^{1-\sigma} \phi_i^{\sigma} \right)^{1/\sigma}}_{\F(\boldsymbol{\phi}, \ba_s)} \right]^{1-\gamma} \right)^{\frac{1}{1-\gamma}}}_{\text{Certainty equivalent over states}}.
\]
This is a nested CES structure: the inner CES aggregates capital types for each state, and the outer CES aggregates states.
\end{calloutimportant}

\paragraph{Special Case: Log Utility ($\gamma = 1$).}
With log utility, the certainty equivalent is the geometric mean:
\begin{equation}
\mathcal{R}(\boldsymbol{\phi}) = \exp\left( \sum_{s=1}^{S} \pi_s \log \F(\boldsymbol{\phi}, \ba_s) \right) = \prod_{s=1}^{S} \F(\boldsymbol{\phi}, \ba_s)^{\pi_s}.
\label{eq:ac-risk-R-log}
\end{equation}

This is the Kelly criterion applied to capital allocation: maximize the expected log return.

\paragraph{Special Case: Deterministic ($S = 1$ or $a_{is} = a_i$ for all $s$).}
With no uncertainty, $\mathcal{R}(\boldsymbol{\phi}) = \F(\boldsymbol{\phi}, \ba)$, and we recover the deterministic portfolio problem from Section~\ref{sec:ac-heterogeneity}.

% --------------------------------------------------------------
\subsection{The Isomorphism: TFP Shocks and Preference Shocks}
\label{subsec:ac-risk-isomorphism}
% --------------------------------------------------------------

A striking feature of the CES framework is the isomorphism between productivity shocks and preference shocks.

\paragraph{Productivity Shocks.}
Consider the production function with stochastic productivity:
\begin{equation}
\F(\bk, \ba_s) = \left( \sum_{i=1}^{N} (\alpha_i a_{is})^{1-\sigma} k_i^{\sigma} \right)^{1/\sigma}.
\label{eq:ac-risk-F-tfp}
\end{equation}

The productivity shock $a_{is}$ multiplies the weight $\alpha_i$, effectively making the ``importance'' of capital type $i$ state-dependent.

\paragraph{Preference Shocks.}
Now consider a certainty equivalent with state-dependent weights:
\begin{equation}
\CE(\bc, \tilde{\boldsymbol{\pi}}_s) = \left( \sum_{i=1}^{N} \tilde{\pi}_{is} \, c_i^{1-\gamma} \right)^{\frac{1}{1-\gamma}},
\label{eq:ac-risk-CE-pref}
\end{equation}
where $\tilde{\pi}_{is}$ is the ``preference weight'' on good $i$ in state $s$.

\paragraph{The Isomorphism.}
Compare the two structures:
\begin{center}
\begin{tabular}{@{}lcc@{}}
\toprule
& \textbf{Production with TFP Shocks} & \textbf{Certainty Equivalent with Pref. Shocks} \\
\midrule
Aggregator & $\F(\bk, \ba_s)$ & $\CE(\bc, \tilde{\boldsymbol{\pi}}_s)$ \\
Power & $\sigma$ & $1 - \gamma$ \\
Weights & $(\alpha_i a_{is})^{1-\sigma}$ & $\tilde{\pi}_{is}$ \\
Inputs & $k_i$ & $c_i$ \\
\bottomrule
\end{tabular}
\end{center}

Setting $\tilde{\pi}_{is} = (\alpha_i a_{is})^{1-\sigma}$ and $1 - \gamma = \sigma$, the two are mathematically identical.

\begin{calloutnote}[Interpretation]
This isomorphism means:
\begin{itemize}[nosep]
\item A shock that raises the productivity of capital type $i$ is equivalent to a shock that raises the ``preference'' for consumption from capital type $i$.
\item The curvature parameter $\sigma$ in production plays the same role as the risk parameter $1-\gamma$ in preferences.
\item Diversification in the production function (spreading investment across capital types) is mathematically equivalent to diversification in a portfolio (spreading consumption across states).
\end{itemize}

This unification is a key insight of the CES framework: \emph{technological complementarity and risk aversion are two sides of the same coin}.
\end{calloutnote}

% --------------------------------------------------------------
\subsection{Ex-Ante Uncertainty: State-Contingent Investment}
\label{subsec:ac-risk-exante}
% --------------------------------------------------------------

Now suppose the state $s$ is revealed \emph{before} investment decisions, or equivalently, the household can choose state-contingent investment $\bI_s$.

\paragraph{Complete Markets.}
With complete Arrow-Debreu markets, the household can trade claims to consumption in each state. Let $q_s$ be the price of one unit of consumption in state $s$.

The household's problem is:
\begin{equation}
\max_{\{C_0, C_{1s}\}} \; u(C_0) + \beta \, \CE(\{C_{1s}\}) \quad \text{s.t.} \quad C_0 + \sum_{s=1}^{S} q_s C_{1s} = W,
\label{eq:ac-risk-complete}
\end{equation}
where $W$ is lifetime wealth.

This is the canonical portfolio problem from earlier chapters: allocate wealth across states to maximize a certainty equivalent, subject to a budget constraint with state prices.

\paragraph{State-Contingent Adjustment Costs.}
Now introduce adjustment costs that apply \emph{after} the state is revealed. In each state $s$, output $Y_s$ is transformed into consumption $C_{1s}$ and investment $I_{1s}$:
\begin{equation}
\Hfun(C_{1s}, I_{1s}) = Y_s = \F(\bk, \ba_s).
\label{eq:ac-risk-H-state}
\end{equation}

The household first observes $s$, then chooses $(C_{1s}, I_{1s})$ subject to the state-specific transformation.

\paragraph{Tobin's $Q$ Becomes State-Dependent.}
In state $s$:
\begin{equation}
Q_s = \frac{\omega_I^{1-\psi}}{\omega_C^{1-\psi}} \left( \frac{C_{1s}}{I_{1s}} \right)^{1-\psi}.
\label{eq:ac-risk-Qs}
\end{equation}

The shadow price of investment varies across states because the consumption-investment ratio varies.

\paragraph{The Euler Equation.}
With state-contingent choices, the Euler equation holds state-by-state:
\begin{equation}
u'(C_0) = \beta \sum_{s=1}^{S} \pi_s \frac{A_s}{Q_s} u'(C_{1s}) \cdot \frac{\partial \CE}{\partial C_{1s}} / u'(\CE),
\label{eq:ac-risk-euler-states}
\end{equation}
where the expression depends on the specific preference structure (expected utility vs. Epstein-Zin).

% --------------------------------------------------------------
\subsection{Comparing Ex-Post and Ex-Ante Uncertainty}
\label{subsec:ac-risk-comparison}
% --------------------------------------------------------------

The two timing conventions lead to fundamentally different problems:

\begin{table}[H]
\centering
\caption{Ex-Post vs. Ex-Ante Uncertainty}
\label{tab:ac-risk-comparison}
\renewcommand{\arraystretch}{1.5}
\begin{tabular}{@{}lcc@{}}
\toprule
\textbf{Feature} & \textbf{Ex-Post} & \textbf{Ex-Ante} \\
\midrule
Information at investment & None (invest, then learn) & Full (learn, then invest) \\
Investment choice & $\bI$ (same in all states) & $\bI_s$ (state-contingent) \\
Adjustment cost applies & Before uncertainty resolves & After uncertainty resolves \\
Tobin's $Q$ & Single $Q$ & State-dependent $Q_s$ \\
CES aggregation & Over states, then over types & Over types, then over states \\
Portfolio problem & Choose $\boldsymbol{\phi}$ to maximize $\CE[\F(\boldsymbol{\phi}, \ba)]$ & Choose $C_{1s}$ to maximize $\CE(\{C_{1s}\})$ \\
\bottomrule
\end{tabular}
\end{table}

\paragraph{The Order of Aggregation.}
The key difference is the \emph{order} in which CES aggregations occur:

\medskip
\noindent
\textit{Ex-post}: First aggregate capital types (production $\F$), then aggregate states (certainty equivalent $\CE$):
\[
\mathcal{R} = \CE_s\left[ \F(\boldsymbol{\phi}, \ba_s) \right].
\]

\medskip
\noindent
\textit{Ex-ante}: First aggregate states (via state prices or certainty equivalent), then aggregate capital types (via transformation):
\[
\text{Welfare} = u(C_0) + \beta \, u\left( \CE_s[C_{1s}] \right), \quad \text{where each } C_{1s} \text{ satisfies } \Hfun(C_{1s}, I_{1s}) = Y_s.
\]

\begin{calloutimportant}[Aggregation Order Matters]
In general, the order of CES aggregations matters---nested CES functions do not commute. This means:
\begin{itemize}[nosep]
\item The optimal investment portfolio differs between ex-post and ex-ante timing.
\item The welfare implications of uncertainty differ.
\item Empirically, identifying the ``true'' timing is crucial for policy analysis.
\end{itemize}

The exception is when all CES aggregators have the same power parameter ($\sigma = 1 - \gamma = \psi$). In this knife-edge case, the aggregations commute and timing is irrelevant.
\end{calloutimportant}

% --------------------------------------------------------------
\subsection{Risk-Adjusted Returns and the Stochastic Discount Factor}
\label{subsec:ac-risk-sdf}
% --------------------------------------------------------------

\paragraph{The Stochastic Discount Factor.}
In the ex-post case, define the stochastic discount factor (SDF):
\begin{equation}
M_s = \beta \frac{u'(C_{1s})}{u'(C_0)}.
\label{eq:ac-risk-sdf}
\end{equation}

With CRRA utility:
\begin{equation}
M_s = \beta \left( \frac{C_{1s}}{C_0} \right)^{-1/\theta}.
\label{eq:ac-risk-sdf-crra}
\end{equation}

\paragraph{Asset Pricing with Adjustment Costs.}
The Euler equation for investment in capital type $i$ is:
\begin{equation}
Q_i = \mathbb{E}\left[ M_s \cdot \frac{\partial \F}{\partial k_i}(\bI, \ba_s) \right],
\label{eq:ac-risk-asset-price}
\end{equation}
where $Q_i$ is the shadow price (Tobin's $Q$) for capital type $i$.

The price of capital equals the expected discounted marginal product, where the expectation is taken under the risk-neutral measure induced by $M_s$.

\paragraph{Risk Premia.}
Capital type $i$ commands a risk premium if its marginal product covaries negatively with consumption:
\begin{equation}
\text{Cov}\left( M_s, \frac{\partial \F}{\partial k_i} \right) < 0 \quad \Longrightarrow \quad Q_i < \mathbb{E}\left[ \frac{\partial \F}{\partial k_i} \right] / R^f,
\label{eq:ac-risk-premium}
\end{equation}
where $R^f = 1 / \mathbb{E}[M_s]$ is the risk-free rate.

\begin{calloutnote}[Adjustment Costs and Risk Premia]
Adjustment costs affect risk premia through two channels:
\begin{enumerate}[nosep]
\item \textbf{Direct}: The transformation elasticity $\eta$ affects how consumption responds to productivity shocks, changing the covariance $\text{Cov}(M_s, \partial \F / \partial k_i)$.
\item \textbf{Indirect}: Adjustment costs slow the reallocation of capital, making the economy less able to respond to shocks, which amplifies risk.
\end{enumerate}

Lower $\eta$ (higher adjustment costs) generally increases risk premia because the economy is ``stuck'' with suboptimal capital allocations when bad states occur.
\end{calloutnote}

% --------------------------------------------------------------
\subsection{The Full Model: Adjustment Costs, Portfolio Choice, and Consumption-Savings}
\label{subsec:ac-risk-full}
% --------------------------------------------------------------

We now assemble the complete model integrating all three elements.

\paragraph{The Planner's Problem (Ex-Post Timing).}
\begin{equation}
\max_{C_0, \bI} \; u(C_0) + \beta \, u\left( \CE\left[ \F(\bI, \ba_s) \right] \right)
\label{eq:ac-risk-full-obj}
\end{equation}
subject to:
\begin{equation}
\Hfun(C_0, \bI) = Y_0.
\label{eq:ac-risk-full-const}
\end{equation}

\paragraph{Decomposition.}
The problem decomposes into:

\medskip
\noindent
\textbf{Step 1: Portfolio choice.} Given total investment $\bar{I}$, choose $\boldsymbol{\phi}$ to maximize:
\begin{equation}
\mathcal{R}^*(\bar{I}) = \max_{\boldsymbol{\phi}: \sum_i \phi_i = 1} \CE\left[ \F(\boldsymbol{\phi}, \ba_s) \right] \cdot \bar{I}.
\label{eq:ac-risk-step1}
\end{equation}

Let $\boldsymbol{\phi}^*$ denote the optimal portfolio (which may depend on $\bar{I}$ if $\Hfun$ induces different effective costs across capital types).

\medskip
\noindent
\textbf{Step 2: Consumption-savings.} Choose $(C_0, \bar{I})$ to maximize:
\begin{equation}
u(C_0) + \beta \, u\left( \mathcal{R}^* \cdot \bar{I} \right) \quad \text{s.t.} \quad \tilde{\Hfun}(C_0, \bar{I}) = Y_0,
\label{eq:ac-risk-step2}
\end{equation}
where $\tilde{\Hfun}$ is the effective transformation given optimal portfolio allocation.

\paragraph{Euler Equation.}
The Euler equation for aggregate investment is:
\begin{equation}
\frac{u'(C_0)}{Q} = \beta \, \mathcal{R}^* \, u'(\mathcal{R}^* \bar{I}) \cdot \frac{\partial \CE}{\partial \text{argument}},
\label{eq:ac-risk-euler-full}
\end{equation}
where the last term accounts for the curvature of the certainty equivalent.

With CRRA utility and $\CE$ as a power mean:
\begin{equation}
u'(C_0) = \beta \frac{\mathcal{R}^*}{Q} \left( \frac{\CE}{\mathcal{R}^* \bar{I}} \right)^{-1/\theta} \CE^{-\gamma} \mathbb{E}\left[ C_{1s}^{-\gamma} \right]^{-1} u'(\CE).
\label{eq:ac-risk-euler-crra}
\end{equation}

The expression simplifies considerably when $\theta = 1/\gamma$ (expected utility) or when $\theta = 1$ (unit IES).

% --------------------------------------------------------------
\subsection{Special Case: Symmetric Shocks and Separation}
\label{subsec:ac-risk-symmetric}
% --------------------------------------------------------------

\paragraph{Symmetric Case.}
Suppose shocks are symmetric across capital types: $a_{is} = a_s$ for all $i$. Then:
\begin{equation}
\F(\bI, \ba_s) = a_s \cdot \F(\bI, \mathbf{1}) = a_s \cdot \left( \sum_{i} \alpha_i^{1-\sigma} I_i^{\sigma} \right)^{1/\sigma}.
\label{eq:ac-risk-symmetric}
\end{equation}

The productivity shock $a_s$ factors out of the production function.

\paragraph{Portfolio Separation.}
With symmetric shocks:
\begin{equation}
\mathcal{R}(\boldsymbol{\phi}) = \CE[a_s] \cdot \F(\boldsymbol{\phi}, \mathbf{1}).
\label{eq:ac-risk-R-symmetric}
\end{equation}

The certainty equivalent of shocks $\CE[a_s]$ is independent of the portfolio. The optimal portfolio maximizes $\F(\boldsymbol{\phi}, \mathbf{1})$, which gives $\phi_i^* = \alpha_i$ as in the deterministic case.

\begin{callouttip}[Separation with Symmetric Shocks]
When productivity shocks are symmetric across capital types:
\begin{itemize}[nosep]
\item The optimal portfolio $\phi_i^* = \alpha_i$ is independent of risk aversion $\gamma$.
\item Risk affects only the consumption-savings decision (through $\CE[a_s]$), not the portfolio.
\item This separation fails when shocks are asymmetric: capital-specific risk induces portfolio tilts.
\end{itemize}
\end{callouttip}

% --------------------------------------------------------------
\subsection{Summary}
\label{subsec:ac-risk-summary}
% --------------------------------------------------------------

\begin{table}[H]
\centering
\caption{Risk and Adjustment Costs: Key Objects}
\label{tab:ac-risk-summary}
\renewcommand{\arraystretch}{1.6}
\begin{tabular}{@{}ll@{}}
\toprule
\textbf{Object} & \textbf{Formula} \\
\midrule
Output in state $s$ & $Y_s = \F(\bk, \ba_s) = \left( \sum_i (\alpha_i a_{is})^{1-\sigma} k_i^{\sigma} \right)^{1/\sigma}$ \\[6pt]
Certainty equivalent & $\CE(C_1) = \left( \sum_s \pi_s C_{1s}^{1-\gamma} \right)^{1/(1-\gamma)}$ \\[6pt]
Risk-adjusted return & $\mathcal{R}(\boldsymbol{\phi}) = \CE_s[\F(\boldsymbol{\phi}, \ba_s)]$ \\[6pt]
Stochastic discount factor & $M_s = \beta (C_{1s}/C_0)^{-1/\theta}$ \\[6pt]
Asset pricing & $Q_i = \mathbb{E}[M_s \cdot \partial \F / \partial k_i]$ \\[6pt]
\bottomrule
\end{tabular}
\end{table}

\begin{callouttip}[Key Insights]
Introducing risk into the adjustment cost framework reveals:
\begin{itemize}[nosep]
\item \textbf{Timing matters}: Ex-post uncertainty (invest then learn) differs fundamentally from ex-ante (learn then invest).
\item \textbf{CES of CES}: The risk-adjusted return $\mathcal{R}$ is a certainty equivalent of a production function---nested CES structures.
\item \textbf{TFP-preference isomorphism}: Productivity shocks and preference shocks are mathematically equivalent in the CES framework.
\item \textbf{Aggregation order}: The order of CES aggregations (over types vs. over states) generally does not commute.
\item \textbf{Risk premia}: Adjustment costs amplify risk premia by limiting the economy's ability to reallocate capital.
\item \textbf{Separation}: With symmetric shocks, portfolio choice separates from risk aversion; with asymmetric shocks, it does not.
\end{itemize}
\end{callouttip}
